\documentclass{article}

\usepackage{mypckg2}

\begin{document}
	
	\begin{itemize}
		\item \textbf{class Loader}. Input: paths of datasets, creates list of datasets (dss) and list of bounds the covered rectangle (bounds) 
		\begin{itemize}
			\item \textbf{lookup}. Input: coordinates of one point, Output: value of the nearest point of the dataset (in our case values from 1 to 20 encoding landscape geodata)
			\item \textbf{lookup\_majority\_window}. Input: coordinates of one point, size of the patch, maximal number of samples, Output: Median value of point laying in the patch around the given point
			\item \textbf{lookup\_list\_majority\_batch\_flat}. Input: list of coordinates, patch size, maximal number of samples, Output: flat list of corresponding values
			\item \textbf{lookup\_list\_majority\_batch\_nested}. Same as lookup\_list\_majority\_batch\_flat but puts out a nested list values.
		\end{itemize}
		\item \textbf{class StitchCoordinates}. Input: stitch length, stitch height, stitch setback, globe diameter, calculates number of initial stitches and number of rows
		\begin{itemize}
			\item \textbf{calculate}: Calculates the number of stitches each row and other quantities
			\item \textbf{coordinates}: Calculates the coordinates of each stitch in each row
			\item \textbf{doublestitches}: Calculates which stitch in each row is a increase, decrease or single crochet. Returns list of lists with entries in \(\{-1,0,1\}\).
		\end{itemize}
		\item \textbf{class PatternGenerator}. Inputs: instance of Loader, instance of StitchCoordinates, Calculates patch size from stitch size, calls \textit{lookup\_list\_majority\_batch\_nested} and stores into \textit{farb}. converts information about stitch type and color into one nested list \textit{self.info}.
		\begin{itemize}
			\item \textbf{info\_globe}. converts information about stitch type and color of the whole globe into one nested list.
			\item \textbf{info\_southern\_hemisphere}. converts information about stitch type and color of the southern hemisphere into one nested list.
			\item \textbf{info\_northern\_hemisphere}. converts information about stitch type and color of the northern hemisphere into one nested list. Differs from the result of the whole globe, because when crocheting the northern hemisphere alone, one starts from the north pole instead of the equator.
			\item \textbf{statistik}. Returns a dictionary, where each color is assigned to the amount of stitches in that color.
		\end{itemize}
		\item \textbf{colorword}. Converts an integer in \(\{1,\ldots,20\}\) representing the landscape type to a name of a color in type of a string.
		\item \textbf{clean\_list}. Converts a list in the format of the info-methods into a compact form, handling repetitions.
	\end{itemize}
	
	% #def calculatenumberofstitches(self):
	%# nos=[self.initialstitches]
	%#for n in range(1, self.numberofrows):
	%# if n < self.equator+ 1:  # southern hemisphere
	%#    if round((math.pi * self.diametermm * math.sin(2 * (n + 1) * self.stitchheight / self.diametermm)+self.stitchsetback) / self.stitchlength) > nos[n - 1] + nos[0] and n!=0:  # no triple increases in the second row
%	#      nos.append(nos[n - 1] + nos[0])
%	#   else:
%	#       nos.append(round(
%	#         (math.pi * self.diametermm * math.sin(2 * (n + 1) * self.stitchheight / self.diametermm)+self.stitchsetback) / self.stitchlength))
%	#else: #northern hemisphere
%	#    nos.append(nos[self.numberofrows-n-1])
%	#self.numberofstitches= nos
	
\end{document}